% !TITLE 行路难·其一
% !AUTHOR 李白
% !DYNASTY 唐
% !ORDER 9

\begin{poem}
金樽清酒斗十千，玉盘珍羞直万钱。
停杯投箸不能食，拔剑四顾心茫然。
欲渡黄河冰塞川，将登太行雪满山。
闲来垂钓碧溪上，忽复乘舟梦日边。
行路难，行路难，多歧路，今安在？
长风破浪会有时，直挂云帆济沧海。
\end{poem}

\begin{appreciation}
这是李白离开长安后写的诗。天宝三载（744年），他在长安待了不到三年，就被赐金放还。那是一场盛宴之后的散场，仰天大笑出门去时的豪情，此刻化作了停杯投箸的茫然。

金樽清酒斗十千，玉盘珍羞直万钱——酒杯是金的，美酒一斗值十千钱；盘子是玉的，珍馐美味值万钱。这是极尽奢华的宴会场景。但停杯投箸不能食，拔剑四顾心茫然——他放下杯子，扔掉筷子，吃不下；拔出剑来，四下张望，心中一片茫然。箸就是筷子。这种反差极具张力：越是物质的丰盛，越衬出精神的空虚。李白不是不喜欢美酒佳肴，而是在这些东西面前，他忽然意识到自己的人生走入了死胡同。

欲渡黄河冰塞川，将登太行雪满山——想渡黄河，冰堵住了河道；想登太行，雪盖满了山头。这是比喻，也是写实。黄河和太行是通往理想的两条路，但都被冰雪封死了。冰和雪是冷的，而他的内心也是冷的。这种内外一致的寒冷，让失望显得更加真实。

闲来垂钓碧溪上，忽复乘舟梦日边——这两句忽然转到了典故。姜太公在渭水边钓鱼，遇周文王而被重用；伊尹梦见自己乘船从日月旁边经过，后来果然被商汤重用。李白用这两个典故告诉自己：暂时的失意不算什么，历史上的伟人也有过困顿的时候。但这种自我安慰是脆弱的，因为紧接着就是一声声的追问：行路难，行路难，多歧路，今安在？

然而，诗的结尾却又昂扬了起来。长风破浪会有时，直挂云帆济沧海——总会有乘风破浪的那一天，那时我将高挂云帆，横渡沧海。这个会有时三个字，透露出一种不确定的坚信：我不知道是什么时候，但我相信一定会有。这种在绝望中硬撑起来的乐观，比一帆风顺时的豪言壮语更打动人。因为它不是无知者的无畏，而是饱经挫折之后的自我激励。
\end{appreciation}
